% Revize 2026-10-27: cílené doplnění; nové úlohy a řešení jsou vlastní tvorba.
% Metodický vzor: 4_rocnik/pripravy/Uvod_do_kombinatoriky.pdf (samostatný dokument).
% Hodina 31, 1. ročník SPŠ. Rozsah: Rozpis_vyuky_MAT_1_rocnik.tex, řádek 31.

% Plán: 1_rocnik/tematicke_plany/SPS.pdf; společná projektová šablona.

% Učebnice: ucebnice/01_Zakladni_poznatky/Mocniny_a_odmocniny/Mocniny a Odmocniny.tex

% Původ úloh: P1: učebnice, příklad za definicí racionální mocniny. Ostatní vlastní. Realisticky.cz: návaznost odmocnin a pravidel pro mocniny, nikoli převzetí definice.

% Doplněný metodický zdroj: externi_zdroje/realisticky/Mocniny/15_MocninySRacionalnimMocnitelem.pdf
% Realisticky.cz 2.7.15: prostudována motivace přechodu k racionálním exponentům,
% převody na odmocniny a gradace příkladů. Příprava zůstává v rozsahu našeho rozpisu.
% URL: http://www.realisticky.cz/ucebnice/01%20Matematika%20S%C5%A0/02%20Funkce%20a%20rovnice/07%20Mocnin%C3%A9%20funkce%20a%20odmocniny/15%20Mocniny%20s%20racion%C3%A1ln%C3%ADm%20mocnitelem.pdf
% Metodický zdroj: externi_zdroje/realisticky/Odmocniny/08_Treti_odmocnina.pdf

% URL: http://www.realisticky.cz/ucebnice/01%20Matematika%20S%C5%A0/01%20Z%C3%A1kladn%C3%AD%20poznatky/02%20%C4%8C%C3%ADseln%C3%A9%20obory/08%20T%C5%99et%C3%AD%20odmocnina.pdf

% Metodický zdroj: externi_zdroje/realisticky/Mocniny/06_Mocniny_s_celym_mocnitelem_II.pdf

% URL: http://www.realisticky.cz/ucebnice/01%20Matematika%20S%C5%A0/01%20Z%C3%A1kladn%C3%AD%20poznatky/06%20Mocniny/06%20Mocniny%20s%20cel%C3%BDm%20mocnitelem%20II.pdf

\documentclass[11pt,a4paper]{article}

% Spolecny styl pro ucitelske pripravy MAT (Work/MAT)
% Verze 1.1, 2026-10-23
\NeedsTeXFormat{LaTeX2e}
\ProvidesPackage{priprava}[2026/10/23 v1.1 Sablona ucitelskych priprav MAT]

% Zakladni typografie a cestina
\RequirePackage[T1]{fontenc}
\RequirePackage[utf8]{inputenc}
\RequirePackage[czech]{babel}
\RequirePackage{lmodern}
\RequirePackage{microtype}
\RequirePackage{amsmath,amssymb}
\RequirePackage{xcolor}
\RequirePackage{enumitem}
\RequirePackage[most]{tcolorbox}
\RequirePackage{tabularx,booktabs,array}
\RequirePackage{fancyhdr}
\RequirePackage{needspace}
\RequirePackage{geometry}
\RequirePackage[hidelinks]{hyperref}

\geometry{a4paper,top=17mm,bottom=18mm,left=23mm,right=23mm}
\definecolor{notegray}{gray}{0.39}
\setlength{\parindent}{0pt}
\setlength{\parskip}{2.5pt plus 1pt minus 0.5pt}
\setlist[itemize]{topsep=2pt,itemsep=1pt,parsep=0pt,leftmargin=1.45em}
\setlist[enumerate]{topsep=2pt,itemsep=2pt,parsep=0pt,leftmargin=1.65em}
\setcounter{secnumdepth}{0}
\renewcommand{\arraystretch}{1.12}

% Zahlavi: vlevo rocnik | globalni tema, vpravo cislo: nazev hodiny
% Zapati zustava zcela prazdne.
\pagestyle{fancy}
\fancyhf{}
\renewcommand{\headrulewidth}{0.35pt}
\renewcommand{\footrulewidth}{0pt}
\setlength{\headheight}{14pt}
\newcommand{\nastavPripravu}[4]{%
  \fancyhead[L]{\small #1\enspace\textbar\enspace #2}%
  \fancyhead[R]{\small #3: #4}%
}

% Minimalisticke cernobile boxy. Kazda definice zvlast.
\tcbset{commonlesson/.style={enhanced,breakable=false,colback=white,colframe=black,
  colbacktitle=white,coltitle=black,boxrule=0.45pt,
  arc=0.8mm,outer arc=0.8mm,left=3mm,right=3mm,top=1.6mm,bottom=1.6mm,
  toptitle=1mm,bottomtitle=0.7mm,fonttitle=\bfseries,
  attach boxed title to top left={xshift=2.5mm,yshift=-1.15mm},
  boxed title style={colback=white,colframe=white,boxrule=0pt,left=1mm,right=1mm}}}
\newtcolorbox{zadaniBox}[1][]{commonlesson,title={Zadání},#1}
\newtcolorbox{prikladBox}[1][]{commonlesson,title={Příklad},#1}
\newtcolorbox{definiceBox}[1][]{commonlesson,title={Definice},#1}

% Spolecny box pro zadani a jeho reseni, bez mezery mezi castmi.
% Pouziti:
%   \begin{prikladReseniBox}
%     Zadani ulohy.
%     \tcblower
%     \begin{align*} ... \end{align*}
%   \end{prikladReseniBox}
% Spodni cast ma jemny sedy podklad a delici linku; lze v ni sazet
% vice radku matematiky. Obe casti maji jen jeden spolecny ramecek.
\tcbset{solutionlesson/.style={
  commonlesson,
  enhanced jigsaw,
  breakable,
  colbacklower=black!3,
  segmentation style={solid,draw=black!25,line width=0.35pt},
  fontlower=\small\color{notegray},
  before lower={\noindent\textbf{Řešení:}\enspace},
  leftlower=3mm,rightlower=3mm
}}
\newtcolorbox{prikladReseniBox}[1][]{solutionlesson,title={Příklad},#1}
\newtcolorbox{zadaniReseniBox}[1][]{solutionlesson,title={Zadání},#1}

% Pomucky pro autora pripravy: sedive poznamky; \reseni pro starsi pripravy
\newcommand{\poznamka}[1]{\par\begingroup\small\color{notegray}\itshape
  \textnormal{\textbf{Poznámka pro učitele.}}\enspace #1\par\endgroup}
\newcommand{\reseni}[1]{\par\begingroup\small\color{notegray}
  \textbf{Očekávané řešení:}\enspace #1\par\endgroup}
\newcommand{\blok}[1]{\needspace{5\baselineskip}\subsection*{#1}}
\newcommand{\vyrok}[1]{\uv{#1}}


\nastavPripravu{1. ročník}{Mocniny}{31}{Mocniny s racionálním exponentem}

\begin{document}

Potřebujeme rozumět zápisu $8^{2/3}$ tak, aby navazoval na mocniny i odmocniny. Nejde o výpočet $8^2:3$. Zlomek v exponentu nejprve přečteme jako jeden celek.

\begin{definiceBox}[title={Definice},breakable=false]
Pro $a\ge0$ a celé $n\ge2$ je \textbf{$n$-tá odmocnina} $\sqrt[n]a$ jediné nezáporné číslo $b$ splňující $b^n=a$. Pro liché $n$ ji lze rozšířit i na záporná $a$, kde má stejný význam vztah $b^n=a$.
\end{definiceBox}


\begin{definiceBox}[title={Definice},breakable=false]
Pro $a>0$, $m\in\mathbb Z$, $n\in\mathbb N$, $n\ge2$ definujeme \textbf{mocninu s racionálním exponentem} předpisem $a^{m/n}=(\sqrt[n]a)^m$. Pro celočíselný exponent platí dřívější definice.
\end{definiceBox}


Čitatel exponentu určuje mocnění, jmenovatel odmocňování. Pro kladný základ lze psát také $a^{m/n}=\sqrt[n]{a^m}$ a platí známá pravidla pro mocniny. Zlomek v exponentu lze krátit. V těchto pravidlech pracujeme s $a>0$; nulu lze doplnit při kladném exponentu, ale záporné základy tímto obecným předpisem nezahrnujeme.

\begin{zadaniReseniBox}[title={Mezikrok: co říká čitatel a jmenovatel exponentu?},breakable=false]
\begin{enumerate}[label=\alph*)]
\item Zapiš pomocí odmocniny $a^{2/5}$ a $a^{-3/4}$ pro $a>0$.
\item Vypočítej $32^{2/5}$, $81^{3/4}$ a $16^{-1/2}$.
\item Uprav $x^{1/3}x^{5/6}$ a $\frac{(a^{3/4})^2}{a^{1/2}}$ pro kladné proměnné.
\end{enumerate}
\tcblower
\begin{enumerate}[label=\alph*)]
\item $(\sqrt[5]a)^2$ a $\frac1{(\sqrt[4]a)^3}$. Jmenovatel exponentu určuje odmocninu, záporné znaménko převrácení.
\item $2^2=4$, $3^3=27$, $1/4$.
\item $x^{2/6+5/6}=x^{7/6}$ a $a^{3/2-1/2}=a$.
\end{enumerate}
\end{zadaniReseniBox}

\begin{prikladReseniBox}[title={Příklad 1},breakable=false]
Vypočítej $8^{2/3}$ a $16^{3/4}$.
\tcblower
$8^{2/3}=(\sqrt[3]8)^2=2^2=4$; $16^{3/4}=(\sqrt[4]{16})^3=2^3=8$. Menší čísla získáme, když nejprve odmocníme.
\end{prikladReseniBox}


\begin{prikladReseniBox}[title={Příklad 2},breakable=false]
Pro $a>0$ uprav $a^{3/2}a^{-1/2}$ a vypočítej $27^{-2/3}$.
\tcblower
$a^{3/2-1/2}=a$. Dále $27^{-2/3}=1/(\sqrt[3]{27})^2=1/9$. Kladnost $a$ ponecháme i u upraveného výrazu.
\end{prikladReseniBox}


\blok{Samostatné procvičení}

\begin{zadaniReseniBox}[title={Procvičení},breakable=false]
\begin{enumerate}[label=\alph*)]
\item Zapiš $\sqrt[5]{a^3}$ pomocí racionálního exponentu při $a>0$.
\item Vypočítej $81^{1/4}$ a $25^{-1/2}$.
\item Pro $x>0$ uprav $\frac{x^{5/3}}{x^{2/3}}$.
\end{enumerate}
\tcblower
\begin{enumerate}[label=\alph*)]
\item $a^{3/5}$; jde o pátou odmocninu třetí mocniny.
\item $3$ a $1/5$.
\item $x^{5/3-2/3}=x$; původní podíl je při $x>0$ definován.
\end{enumerate}
\end{zadaniReseniBox}


\needspace{4\baselineskip}
\poznamka{V učebnici je v definici nepřesná podmínka „$a>b$“ s neurčeným $b$; zde je správně $a>0$. Nezaváděj iracionální exponenty ani logaritmy. Pravidla pro součin odmocnin budou samostatně ve 33. hodině.}


\blok{Navazující procvičování ve třídě}



\begin{zadaniReseniBox}[title={Rezerva: propojení pravidel},breakable=false]
\begin{enumerate}[label=\alph*)]
\item Pro $a,b>0$ uprav $\frac{(a^{2/3}b^{-1/2})^3}{a^{1/2}b^{-1}}$.
\item Proč v našem obecném zápisu racionálních mocnin nepočítáme $(-8)^{2/6}$ přímo stejným předpisem jako pro kladné základy?
\end{enumerate}
\tcblower
\begin{enumerate}[label=\alph*)]
\item $a^{2-1/2}b^{-3/2+1}=a^{3/2}b^{-1/2}=\frac{a^{3/2}}{b^{1/2}}$.
\item Předpis je zaveden pro kladné základy. Šestá odmocnina z $-8$ v $\mathbb R$ neexistuje, zatímco po zkrácení exponentu by vznikla třetí odmocnina. Kladný základ tuto nejednoznačnost vylučuje.
\end{enumerate}
\end{zadaniReseniBox}

\blok{Volitelná rezerva}

\begin{prikladReseniBox}[title={Rozšíření procvičení},breakable=false]
Pro $a,b>0$ uprav $\frac{(a^{1/2}b^{-1/3})^6}{a^2b^{-1}}$.
\tcblower
$\frac{a^3b^{-2}}{a^2b^{-1}}=ab^{-1}=\frac ab$. Všechny úpravy probíhají na zadaném oboru $a,b>0$.
\end{prikladReseniBox}



\poznamka{Nový převod exponentu na odmocninu řeš hned po definici a teprve pak přejdi k původním složeným příkladům. Tato první skupina a původní procvičení tvoří základ. Rezervní otázka o záporném základu je vysvětlení omezení definice, nikoli zavádění nové konvence. Domů zadej uvedených pět položek.}

\blok{Domácí procvičení}

\begin{zadaniReseniBox}[title={Domácí úlohy},breakable=false]
\begin{enumerate}[label=\alph*)]
\item Vypočítej $64^{1/3}$, $125^{2/3}$ a $81^{-1/4}$.
\item Zapiš $\sqrt[3]{a^2}$ a $1/\sqrt[5]{a^3}$ jako mocniny pro $a>0$.
\item Uprav $x^{3/4}x^{-1/2}$ pro $x>0$.
\item Uprav $(a^{1/3}b^{-2/3})^3$ pro $a,b>0$.
\item Uprav $\frac{a^{5/6}}{a^{1/3}}$ pro $a>0$.
\end{enumerate}
\tcblower
\begin{enumerate}[label=\alph*)]
\item $4$, $25$, $1/3$.
\item $a^{2/3}$ a $a^{-3/5}$.
\item $x^{1/4}$.
\item $a/b^2$.
\item $a^{1/2}$.
\end{enumerate}
\end{zadaniReseniBox}

\end{document}
